\begin{derivation}[从\eqnrefs{eq:yukawa-full-si-r10}到\eqnrefs{eq:yukawa-distinguishable-M-si-r10}]
从相互作用

\begin{equation*}
\mathcal L_{\rm int}
=-\sqrt{\hbar c}\,\phi
\left(y_A\bar\psi_A\psi_A+y_B\bar\psi_B\psi_B\right)
\end{equation*}
出发。二阶 Dyson 项含两个 Yukawa 顶角。每个顶角内的费米子双线性型都保持种类，因此外线收缩必须把入射 $A$ 与出射 $A$ 相连、入射 $B$ 与出射 $B$ 相连。于是两顶角之间的标量收缩携带
\begin{equation*}
q=p_1-p_3=p_4-p_2,
\qquad q^2=t.
\end{equation*}
若交换 $p_3\leftrightarrow p_4$，则末态种类也被交换，得到的是不同的外态标签而非同一物理末态的第二种 Wick 收缩，所以可区分粒子没有 $u$-道干涉。

逐个写约化因子，
\begin{align*}
\ii\mathcal M_t
&=\bigl[\bar u_A(p_3)(-\ii y_A)u_A(p_1)\bigr]
\frac{\ii}{t-m_Y^2c^2+\ii0^+}
\bigl[\bar u_B(p_4)(-\ii y_B)u_B(p_2)\bigr],\\
\mathcal M_t
&=\frac{y_Ay_B}{t-m_Y^2c^2}
[\bar u_A(p_3)u_A(p_1)]
[\bar u_B(p_4)u_B(p_2)].
\end{align*}
两个双线性型各有一阶动量量纲，正好抵消传播子的二阶动量量纲，因此约化振幅无量纲，与全书散射归一化一致。
\end{derivation}

\begin{derivation}[从\eqnrefs{eq:yukawa-identical-M-si-r10}到\eqnrefs{eq:yukawa-interference-invariants-si-r10}]
干涉项来自 $t$、$u$ 两种振幅的交叉乘积。对初态两自旋平均并对末态求和后，利用完备关系

\begin{equation*}
\sum_su(p,s)\bar u(p,s)=\slashed p+mc,
\end{equation*}
可把四个双线性型首尾连接成
\begin{equation*}
\mathcal I_{tu}
=\frac14\Tr\!\left[
(\slashed p_3+mc)(\slashed p_1+mc)
(\slashed p_4+mc)(\slashed p_2+mc)
\right].
\end{equation*}
先按斜线记号的个数展开。含一个或三个 gamma 的迹为零，只剩零、二、四个斜线记号：
\begin{align*}
4\mathcal I_{tu}
={}&\Tr(\slashed p_3\slashed p_1\slashed p_4\slashed p_2)\\
&+m^2c^2\Bigl[
\Tr(\slashed p_3\slashed p_1)
+\Tr(\slashed p_3\slashed p_4)
+\Tr(\slashed p_3\slashed p_2)\\
&\hspace{3.2cm}
+\Tr(\slashed p_1\slashed p_4)
+\Tr(\slashed p_1\slashed p_2)
+\Tr(\slashed p_4\slashed p_2)
\Bigr]
+4m^4c^4.
\end{align*}
使用
\begin{align*}
\Tr(\slashed a\slashed b)&=4a\!\cdot b,\\
\Tr(\slashed a\slashed b\slashed c\slashed d)
&=4\bigl[(a\!\cdot b)(c\!\cdot d)
-(a\!\cdot c)(b\!\cdot d)
+(a\!\cdot d)(b\!\cdot c)\bigr],
\end{align*}
得到所需点积形式。下一步不能直接写“用 Mandelstam 变量整理”，而逐个代入等质量关系
\begin{align*}
p_1\!\cdot p_2&=p_3\!\cdot p_4=\frac{s-2m^2c^2}{2},\\
p_1\!\cdot p_3&=p_2\!\cdot p_4=m^2c^2-\frac t2,\\
p_1\!\cdot p_4&=p_2\!\cdot p_3=m^2c^2-\frac u2,
\end{align*}
并用
\begin{equation*}
s+t+u=4m^2c^2
\quad\Longrightarrow\quad
u=4m^2c^2-s-t.
\end{equation*}
四-gamma 部分先化为
\begin{align*}
&(p_3\!\cdot p_1)(p_4\!\cdot p_2)
-(p_3\!\cdot p_4)(p_1\!\cdot p_2)
+(p_3\!\cdot p_2)(p_1\!\cdot p_4)\\
&=\left(m^2c^2-\frac t2\right)^2
-\left(\frac{s-2m^2c^2}{2}\right)^2
+\left(m^2c^2-\frac u2\right)^2.
\end{align*}
六个二-gamma 项的点积和为
\begin{equation*}
\sum_{i<j}p_i\!\cdot p_j
=s+2m^2c^2,
\end{equation*}
其中最后一步由上述三个 Mandelstam 配对逐项相加得到。把 $u=4m^2c^2-s-t$ 代回并展开平方，所有 $u$ 依赖消去，最终得到
\begin{equation*}
\mathcal I_{tu}
=\frac12\left(4m^2c^2s-4m^2c^2t+st+t^2\right),
\end{equation*}
即目标式。质量维数检查给 $[\mathcal I_{tu}]=({\rm momentum})^4$，与两个传播子分母的总四阶量纲相消。
\end{derivation}





\begin{derivation}[从\eqnrefs{eq:yukawa-distinguishable-M-si-r10}到\eqnrefs{eq:yukawa-distinguishable-M2-si-r10}]
取 $m_A=m_B=m$、$y_A=y_B=y$。由\eqnrefs{eq:yukawa-distinguishable-M-si-r10}，对两个初态自旋各平均 $1/2$，并利用完备关系

\begin{equation*}
 \sum_su(p,s)\bar u(p,s)=\slashed p+mc,
\end{equation*}
得到
\begin{equation*}
 \overline{|\mathcal M_t|^2}
 =\frac{y^4}{(t-m_Y^2c^2)^2}\frac14
 \Tr[(\slashed p_3+mc)(\slashed p_1+mc)]
 \Tr[(\slashed p_4+mc)(\slashed p_2+mc)].
\end{equation*}
第一条迹逐项展开为
\begin{align*}
 \Tr[(\slashed p_3+mc)(\slashed p_1+mc)]
 &=\Tr(\slashed p_3\slashed p_1)+m^2c^2\Tr\id_4\\
 &=4(p_3\!\cdot p_1+m^2c^2),
\end{align*}
其中单个 $\gamma$ 矩阵的迹为零。等质量质量壳条件给
\begin{equation*}
 t=(p_1-p_3)^2=2m^2c^2-2p_1\!\cdot p_3,
\end{equation*}
因此
\begin{equation*}
 p_1\!\cdot p_3=m^2c^2-\frac t2,
\end{equation*}
从而
\begin{equation*}
 \Tr[(\slashed p_3+mc)(\slashed p_1+mc)]
 =2(4m^2c^2-t).
\end{equation*}
第二条费米线对应
\begin{align*}
 \Tr[(\slashed p_4+mc)(\slashed p_2+mc)]
 &=4(p_4\!\cdot p_2+m^2c^2).
\end{align*}
由同一个动量转移
\[
 t=(p_2-p_4)^2=2m^2c^2-2p_2\!\cdot p_4
\]
得到
\[
 p_2\!\cdot p_4=m^2c^2-\frac t2,
\]
因此第二条迹也等于 $2(4m^2c^2-t)$。两条迹相乘并乘初态自旋平均因子 $1/4$ 后，
\begin{align*}
 \overline{|\mathcal M_t|^2}
 &=\frac14\,
 \frac{y^4}{(t-m_Y^2c^2)^2}
 [2(4m^2c^2-t)]^2,\\
 &=y^4\frac{(4m^2c^2-t)^2}{(t-m_Y^2c^2)^2}.
\end{align*}
$t$ 与 $m_Y^2c^2$ 都是动量平方，所以结果无量纲；这与约化散射振幅的归一化一致。
\end{derivation}

\begin{derivation}[从\eqnrefs{eq:yukawa-distinguishable-M2-si-r10}到\eqnrefs{eq:yukawa-distinguishable-angular-si-r10}]

正文\eqnrefs{eq:yukawa-distinguishable-M2-si-r10} 已经给出 Lorentz 不变的自旋平均振幅平方。转到质心系后，用 $t=-2p_*^2(1-\cos\theta)$ 和二体相空间通式消去 Mandelstam 变量，即可把不变结果改写成直接可测的角分布~\eqnrefs{eq:yukawa-distinguishable-angular-si-r10}；这一阶段不再引入新的动力学近似。

质心系取
\begin{equation*}
 p_1^\mu=(E/c,\mathbf p),\qquad
 p_2^\mu=(E/c,-\mathbf p),
\end{equation*}
且弹性散射满足 $|\mathbf p_3|=|\mathbf p_1|=p$。于是
\begin{equation*}
 s=(p_1+p_2)^2=\frac{4E^2}{c^2}=4(m^2c^2+p^2).
\end{equation*}
令 $\theta$ 为 $\mathbf p_1$ 与 $\mathbf p_3$ 的夹角，则
\begin{align*}
 t&=(p_1-p_3)^2\\
  &=-|\mathbf p_1-\mathbf p_3|^2\\
  &=-2p^2(1-\cos\theta).
\end{align*}
可区分二体弹性散射的 SI 主式为
\begin{equation*}
 \frac{\dd\sigma}{\dd\Omega}
 =\frac{\hbar^2}{64\pi^2s}
 \frac{|\mathbf p_f|}{|\mathbf p_i|}
 \overline{|\mathcal M|^2}.
\end{equation*}
这里动量比为一。把\eqnrefs{eq:yukawa-distinguishable-M2-si-r10} 与上述 $t$ 代入，即得目标角分布。截面量纲来自 $\hbar^2/s$，因为 $[s]=({\rm kg\,m\,s^{-1}})^2$，故 $[\hbar^2/s]={\rm m^2}$。
\end{derivation}

\begin{derivation}[从\eqnrefs{eq:yukawa-distinguishable-M-si-r10}到\eqnrefs{eq:yukawa-identical-M-si-r10}]
费米场的 Wick 收缩在交换两条外部费米算符时产生一次奇置换。用外态定义

\begin{equation*}
 |p_1s_1,p_2s_2\rangle
 =b_{s_1}^\dagger(p_1)b_{s_2}^\dagger(p_2)|0\rangle
\end{equation*}
和
\begin{equation*}
 \{b_r(\mathbf p),b_s^\dagger(\mathbf p')\}
 =(2\pi\hbar)^3\delta_{rs}\delta^{(3)}(\mathbf p-\mathbf p')
\end{equation*}
交换末态两个产生算符需要一次反对易，因而第二种收缩相对第一种带负号。分别代入两个标量传播子
\begin{equation*}
 \frac{\ii}{t-m_Y^2c^2+\ii0^+},
 \qquad
 \frac{\ii}{u-m_Y^2c^2+\ii0^+},
\end{equation*}
并使用相同的两个 Yukawa 顶角 $-\ii y$，便得到\eqnrefs{eq:yukawa-identical-M-si-r10}。
\end{derivation}

\begin{derivation}[从\eqnrefs{eq:yukawa-NR-M-si-r10}到\eqnrefs{eq:yukawa-potential-SI-r10}]
协变归一化的正能 Dirac 旋量可写成

\begin{equation*}
 u(p)=\sqrt{\frac{E+mc^2}{c}}
 \begin{pmatrix}
 \xi\\[1mm]
 \dfrac{c\boldsymbol\sigma\cdot\mathbf p}{E+mc^2}\xi
 \end{pmatrix}.
\end{equation*}
当 $|\mathbf p|/(mc)\ll1$ 时，
\begin{equation*}
 E=mc^2+\frac{p^2}{2m}+\order\!\left(\frac{p^4}{m^3c^2}\right),
\end{equation*}
因此大分量给
\begin{equation*}
 \bar u(p')u(p)
 =2mc\,\xi'^\dagger\xi
 +\order\!\left(\frac{p^2}{mc}\right).
\end{equation*}
Mandelstam 变量始终满足严格恒等式
\begin{equation*}
 t=-\mathbf q^2+(q^0)^2.
\end{equation*}
定义静态非相对论振幅为同时取 $|\mathbf p|/(mc)\to0$ 与 $q^0/|\mathbf q|\to0$ 的极限。利用上面的旋量展开，便有
\begin{align*}
 \mathcal M_t^{\rm NR}(\mathbf q)
 &\equiv
 \lim_{\substack{|\mathbf p|/(mc)\to0\\ q^0/|\mathbf q|\to0}}
 \mathcal M_t\\
 &=-\frac{4m^2c^2y^2}{\mathbf q^2+m_Y^2c^2}.
\end{align*}
有限速度与有限能量转移的修正由前式中 $O(p^2/m^2c^2)$ 的旋量余项以及精确分母中的 $(q^0)^2$ 项逐项恢复。

Schr\"odinger 第一 Born 近似与上述约化振幅的归一化关系给出
\begin{equation*}
 \widetilde V_Y(\mathbf q)
 =-\frac{4\pi\alpha_Y\hbar^3c}
 {\mathbf q^2+m_Y^2c^2},
 \qquad \alpha_Y=\frac{y^2}{4\pi}.
\end{equation*}
再使用后续推导中得到的三维 Fourier 积分
\begin{equation*}
 \int\frac{\dd^3 q}{(2\pi\hbar)^3}
 \frac{\ee^{\ii\mathbf q\cdot\mathbf r/\hbar}}
 {\mathbf q^2+m_Y^2c^2}
 =\frac{\ee^{-m_Ycr/\hbar}}{4\pi\hbar^2r},
\end{equation*}
得到
\begin{equation*}
 V_Y(r)=-\alpha_Y\hbar c\frac{\ee^{-m_Ycr/\hbar}}{r}.
\end{equation*}
指数中的无量纲组合是 $m_Ycr/\hbar=r/\lambda_Y$，因此 $\lambda_Y=\hbar/(m_Yc)$ 是媒介粒子的 Compton 作用程。$m_Y\to0$ 恢复 $1/r$ 长程极限，$r\gg\lambda_Y$ 时相互作用指数抑制。
\end{derivation}


